% !PRINTED 48
\weberpage{48}
%<<CONT

\begin{equation}
\Delta_0=u_1-u_0,\quad
\Delta_1=u_2-u_1,\quad
\Delta_2=u_3-u_2,\ldots
\tag{2}
\end{equation}
die Reihe ihrer Differenzen, mit
\begin{equation}
\Delta'_0=\Delta_1-\Delta_0,\quad
\Delta'_1=\Delta_2-\Delta_1,\ldots
\tag{3}
\end{equation}
die Reihe ihrer zweiten Differenzen u.\ s.\ f.

Es ist klar, dass die Reihe (1) vollständig bestimmt ist, wenn ihr erstes Glied und die Reihe ihrer ersten Differenzen gegeben ist; denn es ist:
\[
u_1=u_0+\Delta_0,\quad
u_2=u_0+\Delta_0+\Delta_1,\ldots,
\]
\[
u_m=u_0+\Delta_0+\Delta_1+\cdots+\Delta_{m-1}.
\]
Ebenso ist die Reihe (1) völlig bestimmt, wenn die beiden ersten Glieder und die Reihe ihrer zweiten Differenzen gegeben ist u.\ s.\ f. Die Reihe (1) wird eine \gesperrt{arithmetische Reihe $n$ter Ordnung} genannt, wenn die Reihe ihrer $n$ten Differenzen constant ist, also die Reihe der $(n+1)$ten Differenzen aus lauter Nullen besteht.

Man erhält eine arithmetische Reihe $n$ter Ordnung, wenn man in einer ganzen Function $n$ten Grades $f(x)$ für $x$ die Zahlen $0,1,2,3,\ldots$ einsetzt. Denn setzt man:
\[
\Delta_x=f(x+1)-f(x),
\]
\[
\Delta'_x=\Delta_{(x+1)}-\Delta_x,
\]
\[
\cdots\cdots\cdots\cdots\cdots
\]
\[
\Delta_x^{(n-1)}=\Delta_{(x+1)}^{(n-2)}-\Delta_x^{(n-2)},
\]
so ist $\Delta_x$ vom $(n-1)$ten, $\Delta'_x$ vom $(n-2)$ten Grade in Bezug auf $x$ und also $\Delta_x^{(n-1)}$ constant.

Ist nun
\[
u_0,u_1,u_2,\ldots
\]
eine arithmetische Reihe $n$ter Ordnung, so ist die ganze Reihe vollständig bestimmt, wenn die $n$ Werthe $u_0,u_1,u_2,\ldots,u_{n-1}$ und ausserdem die constante $n$te Differenz gegeben sind. Diese letztere ist aber bestimmt, wenn auch noch das $(n+1)$te Glied $u_n$ gegeben ist. Wir können also den Satz aussprechen:

\gesperrt{Eine arithmetische Reihe $n$ter Ordnung ist vollständig bestimmt, wenn ihre $n+1$ ersten Glieder gegeben sind.}

Da nun eine Function $f(x)$ vom $n$ten Grade gleichfalls durch die willkürlich gegebenen Werthe
\[
f(0),f(1),f(2),\ldots,f(n)
\]
%CONT>>
